\hypertarget{chapter-04-page-095}{}
\subsection{具有变核的奇异积分}
\label{sec:ch4-variable-kernel-singular-integrals}

设 $P(x,D)$ 是具有变系数的齐次微分多项式:
\[
 P(x,D)=\sum_{|\alpha|=m}b_\alpha(x)D^\alpha.
\]
由于对 $f\in\mathcal S$ 有
\[
 D^\alpha f(x)=\int_{\mathbb R^n}(2\pi i\xi)^\alpha
 \widehat f(\xi)e^{2\pi ix\cdot\xi}\,d\xi,
\]
所以
\[
 P(x,D)f(x)=\int_{\mathbb R^n}P(x,2\pi i\xi)
 \widehat f(\xi)e^{2\pi ix\cdot\xi}\,d\xi.
\]
用 \eqref{eq:ch4-lambda-definition} 中的 $\Lambda$, 可以得到类似于 \eqref{eq:ch4-differential-operator-factorization} 的表示:
\begin{equation}
\label{eq:ch4-variable-differential-factorization}
 P(x,D)f=T(\Lambda^mf),
\end{equation}
其中
\begin{equation}
\label{eq:ch4-pseudodifferential-operator}
 Tf(x)=\int_{\mathbb R^n}\sigma(x,\xi)\widehat f(\xi)
 e^{2\pi ix\cdot\xi}\,d\xi,
 \qquad
 \sigma(x,\xi)=\frac{P(x,i\xi)}{|\xi|^m}.
\end{equation}
函数 $\sigma$ 关于变量 $\xi$ 是零次齐次的. 把 $\widehat f$ 的定义代入 \eqref{eq:ch4-pseudodifferential-operator}, 形式上得到
\begin{equation}
\label{eq:ch4-variable-kernel-representation}
 Tf(x)=\int_{\mathbb R^n}K(x,x-y)f(y)\,dy,
\end{equation}
其中
\begin{equation}
\label{eq:ch4-variable-kernel-inverse-fourier}
 K(x,z)=\int_{\mathbb R^n}\sigma(x,\xi)e^{2\pi iz\cdot\xi}\,d\xi.
\end{equation}
换言之, 固定 $x$ 后, $K(x,\cdot)$ 是 $\sigma(x,\cdot)$ 的 Fourier 反变换. 由定理~\ref{thm:ch4-smooth-homogeneous-multiplier-representation}, 对每个 $x$ 都存在常数 $a(x)$ 和在 $S^{n-1}$ 上平均值为零的 $\Omega(x,\cdot)\in C^\infty(S^{n-1})$, 使得
\[
 K(x,z)=a(x)\delta(z)+\operatorname{p.v.}
 \frac{\Omega(x,z')}{|z|^n}.
\]

\hypertarget{chapter-04-page-096}{}
这个例子促使我们研究具有变核的奇异积分:
\begin{equation}
\label{eq:ch4-variable-kernel-singular-integral}
 Tf(x)=\lim_{\varepsilon\to0}\int_{|y|>\varepsilon}
 \frac{\Omega(x,y')}{|y|^n}f(x-y)\,dy.
\end{equation}
我们仍可应用旋转方法.

\begin{Theorem}
\label{thm:ch4-variable-odd-kernel-lp}
设 $\Omega(x,y)$ 关于 $y$ 是零次齐次函数, 且满足:
\begin{enumerate}
\item $\Omega(x,-y)=-\Omega(x,y)$;
\item $\Omega^*(u)=\sup_x|\Omega(x,u)|\in L^1(S^{n-1})$.
\end{enumerate}
则 \eqref{eq:ch4-variable-kernel-singular-integral} 给出的算子 $T$ 在 $L^p(\mathbb R^n)$, $1<p<\infty$, 上有界.
\end{Theorem}

\begin{Proof}
由于 $\Omega$ 是奇函数, 可以像推论~\ref{cor:ch4-odd-kernel-lp} 的证明那样得到
\begin{equation}
\label{eq:ch4-variable-kernel-rotation-formula}
 Tf(x)=\frac\pi2\int_{S^{n-1}}\Omega(x,u)H_uf(x)\,d\sigma(u),
 \qquad f\in\mathcal S.
\end{equation}
因此
\[
 |Tf(x)|\le\frac\pi2\int_{S^{n-1}}\Omega^*(u)|H_uf(x)|\,d\sigma(u),
\]
所需结论立即由命题~\ref{prop:ch4-directional-operator-average} 得到.
\end{Proof}

\begin{Theorem}
\label{thm:ch4-variable-kernel-lq-condition}
若把定理~\ref{thm:ch4-variable-odd-kernel-lp} 的条件 2 替换为: 对某个 $1<q<\infty$,
\begin{equation}
\label{eq:ch4-variable-kernel-uniform-lq}
 \sup_x\left(\int_{S^{n-1}}|\Omega(x,u)|^q\,d\sigma(u)\right)^{1/q}
 =B_q<\infty,
\end{equation}
则 $T$ 在 $L^p(\mathbb R^n)$, $q'\le p<\infty$, 上有界.
\end{Theorem}

\begin{Proof}
在 \eqref{eq:ch4-variable-kernel-rotation-formula} 中应用指数为 $q$ 与 $q'$ 的 Hölder 不等式, 得
\begin{equation}
\label{eq:ch4-variable-kernel-holder-bound}
 |Tf(x)|\le\frac\pi2B_q
 \left(\int_{S^{n-1}}|H_uf(x)|^{q'}\,d\sigma(u)\right)^{1/q'}.
\end{equation}
将两边取 $q'$ 次幂并关于 $x$ 积分, 得到 $p=q'$ 时的估计. 若 \eqref{eq:ch4-variable-kernel-uniform-lq} 对某个 $q$ 成立, 则它对每个更小的指数也成立, 因而得到全部 $q'\le p<\infty$.
\end{Proof}

\hypertarget{chapter-04-page-097}{}
当 $\Omega$ 为偶函数时, 有一个与定理~\ref{thm:ch4-variable-kernel-lq-condition} 类似的结果, 其证明与定理~\ref{thm:ch4-rough-kernel-lp} 的证明相仿. 它见第~\ref{subsec:ch4-references}~节所引 Calderón 与 Zygmund 的第一篇论文.

形如 \eqref{eq:ch4-pseudodifferential-operator} 的算子是一类特殊的伪微分算子. 更多内容见\MBUnresolvedReference{section}{第 5 章第 6.9 节}.
